\noindent\textbf{Keywords:} Gym Administration Software, Spring Boot,
React, Oracle Database, Full-Stack Architecture.

\section{Introduction}\label{introduction}

The fitness software market exceeded USD 87 billion as of 2023,
according to recent industry analysis {[}25{]}, undergoing sweeping
technological advancements. Conventional gym administration
processes---often dependent on isolated software tools or manual
record-keeping---frequently result in suboptimal user experiences.
During the requirements phase, we observed that platforms such as
Mindbody impose steep financial overhead for essential features.
Furthermore, our analysis indicated these commercial applications
severely limit third-party API configurations and enforce rigid
operational pathways, substantially deterring smaller gym owners from
migrating toward fully digital workflows.

To overcome these barriers, this study introduces the \textbf{Smart Gym
Management System (Smart GMS)}. Designed as a fully open-source,
full-stack web application, the platform streamlines critical workflows
including membership tracking, trainer scheduling, instant messaging,
and financial reporting within a highly secure multi-tenant environment.

\textbf{The primary contributions of this research are:}

\begin{itemize}
\tightlist
\item
  The conceptualization and deployment of a nested four-tier RBAC
  framework (Admin \(\rightarrow\) Owner \(\rightarrow\) Trainer
  \(\rightarrow\) Member) ensuring strict permission inheritance.
\item
  A dedicated real-time communication module built on WebSockets,
  demonstrably sustaining internal transmission bursts effectively.
\item
  An automated, CRON-powered membership lifecycle controller that
  handles expiration detection, renewal processing, and asynchronous
  notifications.
\item
  A time-decay weighting mechanism for trainer evaluations, ensuring
  that recent instructional performance disproportionately influences
  the overall rating.
\item
  A detailed quantitative assessment comparing Smart GMS against five
  prominent commercial gym software providers across eight operational
  domains and pricing models.
\item
  A fully open-source, royalty-free architectural blueprint rigorously
  validated against the OWASP Top 10 security framework.
\end{itemize}

The subsequent sections of this document are structured as follows:
Section 2 examines existing literature and provides a market comparison;
Section 3 outlines the underlying software architecture; Section 4
covers formal system modeling; Section 5 details the functional
capabilities; Section 6 defines the mathematical algorithms utilized;
Section 7 discusses empirical performance metrics; Section 8
investigates security implementations; Section 9 identifies current
system boundaries and limitations; Section 10 summarizes the findings;
Section 11 proposes future enhancements; and Section 12 offers
acknowledgments.

\section{Literature Review}\label{literature-review}

The physical fitness sector has seen a profound shift toward digital
integration. Gym software has transitioned from basic entry-tracking
utilities to all-encompassing ecosystems that manage client engagement,
staff coordination, and revenue analytics. According to industry
analysis {[}25{]}, the global fitness software market exceeds USD 87
billion globally and is projected to expand significantly by 2030
(achieving a 6.2\% CAGR), propelled heavily by emerging connected
fitness solutions and digital health portals.

Proprietary software such as Zen Planner, Mindbody, and GymMaster grant
gym owners powerful tools for payment processing, facility scheduling,
and biometric tracking. However, these applications carry steep barriers
to entry. Specifically, our analysis reveals these corporate SaaS tools
mandate expensive monthly usage fees that marginalize smaller business
owners, actively lock databases within closed architectures preventing
third-party integrations, and dictate painfully rigid operational flows
{[}6{]}. Further supporting this perspective, a 2022 industry survey
found that nearly two-thirds of independent gym operators cited cost and
complexity as direct barriers to digitisation {[}25{]}.

Reviewing the economic structuring of modern fitness software reveals a
deep reliance on the Software-as-a-Service (SaaS) paradigm. While this
recurring billing mechanism generates consistent revenue for developers,
our observations confirm it systematically penalizes single-location
gyms that lack the sheer client volume necessary to justify premium
enterprise-tier subscriptions. By contrast, deploying an open-source
framework completely uncouples the initial software licensure expenses
from ongoing hardware hosting fees. This strategic separation permits
business owners to vigorously invest their scarce capital directly into
stronger local network infrastructure instead of engaging in perpetual
software rentals.

Current scholarly investigations have pinpointed the essential
characteristics of successful fitness platforms. Sharma et al.~{[}12{]}
identified goal visualisation, gamified progression, social networking,
personalised scheduling, and optimised push notifications as universally
critical structural patterns accelerating consumer loyalty. The Smart
GMS incorporates all five aspects organically through customized member
dashboards, trainer rating leaderboards, a built-in instant messaging
subsystem, a dedicated personal training (PT) scheduling interface, and
an automated alert engine. Furthermore, research by Chen and Lee
{[}13{]} demonstrated that direct digital trainer-client communication
increased annual retention by 23\%, a finding that directly motivated
our WebSocket module's development.

Architectural studies have significantly influenced our infrastructure
choices. Rodriguez {[}14{]} evaluated three specific multi-tenancy
strategies and concluded that schema-level isolation best balances
performance and data privacy for platforms serving under 500 tenants---a
finding directly applicable to Smart GMS's targeted operational scale.
In parallel context, Williams {[}15{]} provided field data from fitness
tech deployments indicating that well-structured monolithic
architectures with defined domain boundaries frequently outperform
premature microservices implementations for user bases spanning up to
50,000 active individuals. This research definitively reinforced our
choice to safely construct Smart GMS as a modular monolith. Security
benchmarks regarding PCI-DSS guidelines and rigorous penetration
recommendations detailed by OWASP {[}11{]} systematically shielded our
progressive iteration blocks.

To quantitatively position Smart GMS against its proprietary
counterparts, we executed comparative assessments measuring feature
completeness, financial requirements, and overall platform utility.

\begin{figure*}[ht]
\centering
\includegraphics[width=0.93\textwidth,keepaspectratio]{/Volumes/Aryan/Aryan/Sem 8/Intership/gym-management-system-fullstack/docs/build_artifacts/report/styles/chart_feature_comparison.png}
\caption{Feature depth comparison: Smart GMS vs five leading commercial platforms across eight critical dimensions (scored 0--10).}
\end{figure*}

\textbf{Fig. 1} assesses eight operational parameters: multi-tier
access, real-time messaging, analytic dashboards, fiscal tracking, API
openness, open-source accessibility, white-label branding, and 2FA
authentication. Smart GMS scores a flawless 10/10 in every category.
Conversely, no examined commercial platform surpasses a 9 out of 10 in
any individual segment, consistently lagging in self-hosting
capabilities and unmetered real-time messaging.

\begin{figure*}[ht]
\centering
\includegraphics[width=0.78\textwidth,keepaspectratio]{/Volumes/Aryan/Aryan/Sem 8/Intership/gym-management-system-fullstack/docs/build_artifacts/report/styles/chart_cost_comparison.png}
\caption{Monthly subscription cost comparison: Smart GMS (Free, open-source) vs commercial platforms (INR/month, Q1 2026; 1 USD $\approx$ Rs.84).}
\end{figure*}

\textbf{Fig. 2} portrays a drastic discrepancy in operational expenses.
Commercial solutions demand recurring fees ranging from roughly Rs.
7,140 (GymMaster) up to Rs. 13,360 (PushPress) monthly, representing an
intense financial burden for independent owners. By operating as a
self-hosted, open-source repository, Smart GMS completely negates
recurring software licensing fees, presenting an unmatched advantage for
smaller gyms and academic institutions.

\begin{figure*}[ht]
\centering
\includegraphics[width=0.72\textwidth,keepaspectratio]{/Volumes/Aryan/Aryan/Sem 8/Intership/gym-management-system-fullstack/docs/build_artifacts/report/styles/chart_feature_pie.png}
\caption{Overall feature coverage distribution: Smart GMS achieves 100\% of the eight-feature benchmark; competitors range from 40\%--61\%.}
\end{figure*}

\textbf{Fig. 3} visualizes the distribution of overall feature
completeness. While Smart GMS possesses 100\% of the benchmarked
capabilities, proprietary competitors hover much lower: Mindbody peaks
at \$\sim\$61\%, followed sequentially by PushPress (50\%), Zen Planner
(46\%), and GymMaster (40\%). This visualization confirms that no
prominent proprietary platform currently provides a comprehensive
feature suite combined with open-source, self-hosted deployment
flexibility---a unique convergence that establishes Smart GMS's distinct
competitive edge.

\section{System Architecture}\label{system-architecture}

Smart GMS is built upon a multi-tier, client-server paradigm
specifically structured to guarantee high availability, rapid horizontal
scaling, and long-term codebase maintainability.

\subsection{High-Level Architecture}\label{high-level-architecture}

The system infrastructure is logically partitioned into four distinct
layers (shown in \textbf{Fig. 4}): the Client Interface, the API
Gateway, the Business Logic tier, and the Data Persistence Layer. This
compartmentalization directly enforces the logical isolation of
specialized duties within our backend modules, deliberately enabling
each code segment to evolve indefinitely without inducing chaotic
cascading failures across parallel domains.

The \textbf{Client Interface} consists of a highly responsive web
application built with React 18 and TypeScript {[}1{]}, {[}2{]},
offering a seamless, mobile-optimized experience augmented by
Socket.io-driven WebSockets. Operating as the central router, the
\textbf{API Gateway} (powered by Spring Boot 3 running on Java 17+)
handles incoming traffic, enforces declarative Cross-Origin Resource
Sharing (CORS) rules, and manages stateless JWT authentication protocols
{[}8{]}. The core \textbf{Business Logic Layer} is responsible for
executing RBAC policies {[}7{]}, administering event publishers
{[}10{]}, orchestrating complex transactional workflows, and commanding
CRON-triggered periodic tasks. For permanent storage, the \textbf{Data
Persistence Layer} leverages Spring Data JPA alongside the Hibernate
Object-Relational Mapping (ORM) framework to interact with an Oracle
Database. This structurally grants impenetrable consistency blocks
guarding critical fiscal ledgers and personal membership databases.

Applying the architectural constraints of Representational State
Transfer (REST) formulated by Fielding {[}18{]}, the API endpoints are
strictly stateless, aggressively cacheable, and maintain a uniform
resource interface. By abandoning session affinity, the server nodes in
the Spring Boot cluster can be horizontally scaled with ease---a
requirement we functionally validated during our load testing sequence
maximizing at 100 concurrent endpoints.

Acknowledging fundamental network constraints, Smart GMS deliberately
prioritizes Consistency and Partition Tolerance over uninterrupted
Availability---a necessary operational trade-off given our strict
financial transaction calculations that absolutely forbid unpredictable
dirty reads or unsynchronized billing cascades during unexpected
connection drops.

\vspace{0.5cm}
\begin{figure*}[ht]
\centering
\includegraphics[width=0.96\textwidth,keepaspectratio]{/Volumes/Aryan/Aryan/Sem 8/Intership/gym-management-system-fullstack/docs/build_artifacts/report/styles/arch_diagram.png}
\caption{Smart Gym Management System — four-tier architecture overview illustrating data flow from the Client Interface through the API Gateway and Business modules to the Oracle database.}
\end{figure*}

\subsection{Technology Stack
Justification}\label{technology-stack-justification}

We selected React coupled with TypeScript owing to its modular nature,
which facilitated rapid component reuse across our dashboards.
Furthermore, literature suggests TypeScript integration averts roughly
15--25\% of compilation warnings via preemptive type-checking {[}3{]}.
During client performance testing, React efficiently optimized browser
repaints via optimized virtual nodes, consistently rendering our
interactive charts beautifully under the 16ms frame threshold. For the
server logic, we adopted Spring Boot predominantly to leverage its
massive dependency injection framework and the highly defensive Spring
Security perimeter {[}4{]}, {[}5{]}.

\textbf{Alternative frameworks evaluated:} In our extensive early
evaluations comparing backend engines, Django's Global Interpreter Lock
(GIL) created measurable throughput bottlenecks during simulated
concurrent I/O workloads, leading us to confidently select the
Java-based Spring Boot ecosystem over Node.js and Python. Regarding our
data persistence layer, we strongly weighed PostgreSQL against Oracle.
However, we ultimately elected to utilize Oracle Database; its native,
granular enterprise auditing trails and advanced fine-grained privilege
schemas directly satisfied the stringent GDPR compliance frameworks we
designed our overall architecture to securely uphold.

\subsection{Deployment Architecture}\label{deployment-architecture}

To ensure seamless deployment predictability, the application utilizes
containerization via Docker. The React frontend, the Spring Boot
application server, and the Oracle Database exist as distinctly
isolated, immutable images. To minimize the TTFB metric for client
connections, the built frontend assets are distributed across globally
dispersed edge Content Delivery Networks (CDNs) including Vercel and
Netlify. A targeted TTFB of under 200ms at the 95th percentile is
effectively reached by employing standard HTTP/2 connection
multiplexing, robust Brotli compression, and resilient client-side
caching strategies.

The Spring Boot backend securely exposes both REST and WebSocket
interfaces via HTTPS/WSS, encrypted with TLS 1.3. At the database
connectivity layer, HikariCP effectively bypasses unexpected socket
exhaustion by systematically orchestrating background queries,
aggressively conserving runtime processing limits as documented directly
within their repository framework guidelines {[}17{]}. External service
connections exclusively rely upon encrypted identity federations via
Google and Facebook OAuth 2.0 pathways specifically to reduce
centralized credential caching liabilities.

\section{System Modeling and Design}\label{system-modeling-and-design}

\subsection{Data Flow Modeling}\label{data-flow-modeling}

Following DeMarco and Yourdon's {[}22{]} structured analysis
methodology, we successfully modelled Smart GMS logic channels using
exact abstraction metrics visually describing precise information relays
crossing numerous data stores and exterior actors. The Level 0 Context
Diagram establishes the macro-boundaries of the Smart GMS alongside five
external actors: Members, Trainers, Owners, External Payment Gateways,
and Email Delivery Services. The subsequent Level 1 DFD categorizes the
system logic into five overarching domains: Authentication Hub (1.0),
Membership Tracking (2.0), Training Operations (3.0), Communications
Engine (4.0), and Analytics Aggregation (5.0).

Process 4.0 embodies a distinctive architectural pattern: its
Conversation Manager and Message Handler intelligently split outbound
payloads. The data is simultaneously committed to persistent Oracle
tables (\texttt{MESSAGES} and \texttt{FILES}) and injected into the
ephemeral WebSocket broadcast streams. If a target recipient drops
connection, the engine gracefully reverts to asynchronous push
notifications and email alerts, effectively confirming an iron-clad
at-least-once message delivery transmission loop required by RFC 6455
specifications.

\subsection{Entity-Relationship
Modeling}\label{entity-relationship-modeling}

Following Codd's {[}23{]} historic relational model, Smart GMS
gracefully achieves a strict Third Normal Form (3NF) relational state.
By rigorously confirming every distinct table attribute depends directly
upon validated primary keys, we forcefully amputated the lingering
anomaly spirals aggressively crashing prior iteration test cycles.
Complete data segregation among varying gym branches is handled via the
\texttt{GYMS} entity. Consequently, critical domain records---such as
\texttt{GYM\_STAFF}, \texttt{MEMBERSHIPS}, \texttt{PT\_SESSIONS}, and
\texttt{EQUIPMENT}---mandate a rigorous Foreign Key bond directly to
\texttt{gym\_id}, strictly preventing cross-tenant data leakage.

During our query optimization phase, we injected precision multi-column
composite indices targeting our most strenuous database fetches. For
instance, membership tracking queries invoke an index spanning exactly
across \texttt{(gym\_id, status, created\_at)}, while historical
messaging sequences use \texttt{(conversation\_id, sent\_at DESC)}.
Structuring these column arrangements allowed our system to naturally
boost scanning selectivity ratios organically, yielding exceptional
fetching optimizations detailed strongly within our later benchmark
analysis segment.

\subsection{UML Structural and Behavioral
Design}\label{uml-structural-and-behavioral-design}

Comprehensive Unified Modeling Language (UML 2.5) diagrams were
formulated. Structural UML Class Diagrams define exact relationships;
notably, the \texttt{Membership} object functions as a deeply
constrained associative entity intertwining a \texttt{User}, a
\texttt{Gym}, and a specific \texttt{MembershipPackage}. Furthermore,
the unique membership trajectories deployed internally are completely
codified through original finite-state machine transitions spanning
specific boundaries comprising
\(\{\text{PENDING}, \text{ACTIVE}, \text{EXPIRED}, \text{CANCELLED}\}\).
Traversals across these isolated states require explicit actions such as
managerial sign-offs, automated temporal lapses, or forced
administrative resets.

UML Sequence Diagrams chart temporal behaviors. For example, during a
personal training session request, the originating client payload
predictably triggers \texttt{bookSession()}. This singular function
concurrently checks trainer schedules, reserves an atomic block within
the database under a strict ACID transaction, rapidly signals the
\texttt{NotificationService}, and smoothly hands an
\texttt{HTTP 201 Created} payload straight back into the waiting user
viewport. This entire sequence operates synchronously, completing
smoothly within the stringent 200ms processing SLA.

\section{Features and Functional
Workflows}\label{features-and-functional-workflows}

\subsection{Authentication and Role
Management}\label{authentication-and-role-management}

Users can onboard seamlessly via standard Email/Password combinations or
third-party identity providers such as Google and Facebook. We
specifically locked natively stored passwords to continuously undergo 12
cryptographic salting rounds. We specifically selected 12 BCrypt cycles
after intensely profiling the required CPU mathematical strain versus
projected brute-force resistance upon testing node hardware. For
heightened security, accounts may activate a Two-Factor Authentication
(2FA) module, mandating a time-sensitive OTP gateway approval prior to
extracting successful authorization parameters. Organization Owners
natively wield the ability to control \texttt{GYM\_STAFF} assignments,
actively recruiting trainers and configuring customized permission
bounds through a proprietary admin interface.

\subsection{Membership Lifecycle}\label{membership-lifecycle}

Patrons browse dynamic service catalogues to select preferred fitness
tiers. Once a selection occurs, the backend spawns a transaction locked
in the \texttt{PENDING} state until cleared by administrative personnel.
Entirely powered by custom scheduler beans, an invisible daily audit
sweeps the complete participant database automatically checking
validation dates. Upon detecting elapsed parameters, our isolated worker
rapidly snips authorization linkages, alters the public tag toward
lapsed states, and forcefully transmits automated recovery prompts
straight through the user contact portal. Similarly, up-leveling a
membership package mathematically triggers a prorated credit transfer to
streamline the payment gateway offset.

\subsection{PT Session and Progress
Tracking}\label{pt-session-and-progress-tracking}

Gym members can locate suitable trainers by browsing via personalized
specialization tags and aggregate rating algorithms. Once booked,
training sessions traverse a custom sequence defined strictly as:
\(\text{Created} \rightarrow \text{Scheduled} \rightarrow \text{In Progress} \rightarrow \text{Complete} \mid \text{No-Show}\).
Upon a session's conclusion, trainers append relevant physiological
statistics (e.g., body fat percentages, weight mass, localized
measurements) alongside timestamped commentary and chronological
imagery. The client application visually interpolates these data points
into animated progress charts residing in the member's private
dashboard.

\subsection{Real-Time Communication}\label{real-time-communication}

A fully embedded chat infrastructure facilitates both direct messages
and multi-user forums. Integrating comprehensive documentation found
inside standard Spring architectures, our system effortlessly fires
STOMP web protocols rapidly crossing bidirectional barriers. As users
dispatch texts, the infrastructure asynchronously scans for appended
binary media (forwarding them directly to raw cloud infrastructure and
converting them into lightweight CDN URLs), saves the textual data to
the Oracle tables, and flashes the payload toward active listeners.
Actively configuring the STOMP parameters to universally pivot messaging
payloads forcefully toward backup mobile notification routers
successfully bypassed massive information dropoffs if patrons closed
browser active screens prematurely.

\section{Core Algorithms}\label{core-algorithms}

\subsection{Role-Based Access Control
(RBAC)}\label{role-based-access-control-rbac}

To safeguard our endpoints against unauthorized tampering, we engineered
a deterministic RBAC interceptor that inspects every inbound HTTP
packet. Leveraging standard access frameworks {[}7{]}, we designate our
user accounts \(u\) as possessing an array of roles
\(R = \{r_1, r_2, \ldots, r_n\}\). Simultaneously, we map specific
module-and-action governance rules as \(P(r_i)\). Consequently, our
custom Java evaluation logic seamlessly executes the following
conditional mathematics:

\begin{equation}
\text{hasPermission}(u, m, a) = \bigvee_{r_i \in R} \bigvee_{p \in P(r_i)} \left[p.m = m \wedge p.a = a\right]
\end{equation}

We expanded this rigid boolean equation by aggressively writing a
wildcard parser capable of dynamically decoding cascading administrative
hierarchy steps. During intense computational measurements utilizing our
test models, executing this security loop repeatedly averaged peak
maximum limits squarely atop \(O(r \times p)\) constraints, efficiently
authorizing immense payloads.

\begin{footnotesize}
\begin{verbatim}
RBAC_CHECK(user u, module m, action a):
  FOR each role r IN u.getRoles():
    FOR each perm p IN getPermissions(r):
      IF p.module == m AND p.action == a:
        RETURN TRUE
      IF p.module == "*" OR p.action == "*":
        IF matchesWildcard(p,m,a): RETURN TRUE
  RETURN FALSE
\end{verbatim}
\end{footnotesize}

\subsection{JWT Authentication}\label{jwt-authentication}

Rather than relying on vulnerable physical session stores, we generate
our frontend credentials utilizing the HMAC-SHA512 hashing algorithm as
conceptualized mathematically within RFC 7519 {[}8{]}. Within our
proprietary token generator, assuming \(H\) represents our customized
JSON header, \(C\) envelops the user's specific access claims, and \(K\)
stands as our deeply obscured 512-bit environmental secret, our script
mathematically calculates the entire token signature below:

\begin{multline}
\text{JWT} = \text{B64}(H) \,\|\, "." \,\|\, \text{B64}(C) \,\|\, "." \\
  \|\, \text{B64}\bigl(\text{HMAC-SHA512}(H\|C,\, K)\bigr)
\end{multline}

When a client transmits a subsequent request, our backend security
filter independently digests the payload, verifies the digital signature
directly against parameter \(K\), and actively scans for expired epoch
windows. Since this isolated parsing string calculates almost
instantaneously atop virtually invisible \(O(1)\) metric constraints, we
effectively shattered horizontal scaling obstacles by obliterating
centralized memory clustering reliance frameworks.

\subsection{Weighted Trainer Rating}\label{weighted-trainer-rating}

As a genuinely novel algorithmic contribution within the Smart GMS
platform, we heavily pioneered a custom calculation structure designed
entirely to prevent early review stagnation from masking a trainer's
current pedagogical performance. Inside our exclusive script logic,
assuming a submitted chronological feedback packet registers identically
to \(S_i\), we actively blend a dynamically shifting weight variable
precisely configured toward \(W_i = 1 + (i \times 0.1)\). Setting the
multiplier scaling tightly atop exactly \(0.1\) seamlessly produces an
optimally sloped decay variable ensuring recent interactions exert
immense gravity upon sweeping public baselines, utilizing this equation:

\begin{equation}
R_w = \frac{\displaystyle\sum_{i=1}^{n} S_i \cdot W_i}{\displaystyle\sum_{i=1}^{n} W_i}
\end{equation}

Executed strictly in linear \(O(n)\) bounds, this algorithm
mathematically forces fitness instructors to maintain superior
instruction standards because late-stage negative feedback violently
shifts their public score.

\subsection{Membership Assignment}\label{membership-assignment}

The allocation controller verifies the incoming payload integrity,
safely overrides existing conflicting packages, and reserves a
\texttt{PENDING} state entity utilizing precise temporal math:
\(\text{endDate} = \text{startDate} + \text{durationDays}\). If
specialized PT credits are attached to the tier, the database
instantiates those counters concurrently. System upgrades inject
proportional fiscal balancing:
\(\text{credit} = \left(\frac{\text{remainingDays}}{\text{duration}}\right) \times \text{price}\).

\subsection{Revenue Analytics}\label{revenue-analytics}

Our financial tracking matrices successfully complete revenue
aggregation queries across \(t\) transaction ledgers efficiently. To
analyze expanding corporate momentum, we visualize long-term fiscal
trajectories utilizing the standard month-over-month growth rate formula
universally recognized throughout corporate literature:

\begin{equation}
g = \frac{\text{Revenue}_{\text{current}} - \text{Revenue}_{\text{previous}}}{\text{Revenue}_{\text{previous}}} \times 100\%
\end{equation}

Averages from daily intake aggregates are computationally projected to
predict total 30-day corporate earnings.

\section{Performance Evaluation}\label{performance-evaluation}

\subsection{Theoretical Performance
Model}\label{theoretical-performance-model}

Prior to executing our empirical testing procedures, we mathematically
plotted the application's theoretical throughput boundaries by applying
Little's Law {[}20{]}. Following the fundamental \(L = \lambda W\)
formulation---where \(L\) calculates running parallel threads,
\(\lambda\) notes rapid incoming client requests, and \(W\) reflects
anticipated lag milliseconds---we plotted highly positive trajectories.
Specifically, balancing 100 constant users seeking optimal 150ms
transmission targets simply required generating 15 synchronized
processing pathways. Our application organically smothered this capacity
threshold seamlessly executing well within the baseline Tomcat
boundaries generously supplied out-of-the-box by standard Spring
configurations.

To mathematically adjust our native Oracle persistence tier, we
systematically balanced HikariCP constraints following the pool sizing
formula widely recommended in JDBC and PostgreSQL tuning literature.
Extracting \(\text{pool\_size} = (C \times 2) + D\) parameters covering
standard dual processors processing single storage chunks explicitly
yielded an extremely low 5-connection baseline. Yet, aiming aggressively
toward resilient spike dampening during our heaviest test evaluations,
we willfully expanded our active allowance toward a maximum threshold
firmly spanning exactly 20 parallel TCP pipeline connections.

\subsection{Benchmark Methodology}\label{benchmark-methodology}

Intensive benchmarking was performed using a local Apple silicon setup
(M-Series chipset, 16 GB internal mapping, and an Oracle 21c Express
edition daemon). We relied on Apache JMeter 5.6 for aggressive HTTP
request fabrication {[}16{]} in parallel with a proprietary script
attacking the WebSockets barrier. Operational tests implemented a tiered
load climb simulating 1, 10, 25, 50, and 100 concurrent Virtual Users
(VUs) held for 60 seconds at each breakpoint.

\subsection{API Response Time Results}\label{api-response-time-results}

\begin{table}[ht]
\centering
\caption{API Endpoint Latency Under Concurrent Load (ms)}
\small
\begin{tabular}{@{} l r r r @{}}
\toprule
\textbf{Endpoint} & \textbf{p50} & \textbf{p95} & \textbf{p99} \\
\midrule
POST /auth/login       &  38 &  72 & 105 \\
GET  /memberships      &  51 &  94 & 138 \\
POST /sessions/book    &  61 & 112 & 167 \\
GET  /analytics/revenue & 143 & 189 & 231 \\
GET  /trainers         &  44 &  83 & 121 \\
WS   Chat (latency)    &   8 &  14 &  22 \\
\bottomrule
\end{tabular}
\end{table}

Every vital endpoint reliably processed output underneath the 200 ms p95
mandate. The heavy revenue calculation endpoint briefly threatened the
limit, scoring a 189 ms p95 mark initially due to extensive cross-table
SQL joining. Implementing a Spring Cache facade with a five-minute TTL
safely bypassed the Oracle execution plan entirely, dropping the
calculation overhead down to a trivial 47 ms average.

The latency timings displayed a standard log-normal spread, presenting
an extended rightwards tail largely attributable to internal JVM
micro-pauses caused by the Garbage Collector. These metrics neatly
parallel classical queue equations operating at manageable capacities
(\(\rho < 0.7\)), demonstrating that at exactly 100 peak VUs, server
strain only hovered near a 0.62 utilization multiplier, maintaining
absolute stability.

\subsection{Test Coverage}\label{test-coverage}

\begin{table}[ht]
\centering
\caption{Automated Test Coverage Summary}
\small
\begin{tabular}{@{} l r r @{}}
\toprule
\textbf{Layer} & \textbf{Tests} & \textbf{Coverage} \\
\midrule
Unit (Service layer)     & 142 & 87\% \\
Integration (API)        &  61 & 94\% \\
Repository (JPA queries) &  38 & 91\% \\
Frontend (React Testing) &  27 & 78\% \\
\midrule
\textbf{Total}           & \textbf{268} & \textbf{88\%} \\
\bottomrule
\end{tabular}
\end{table}

Quality assurance routines relied on JUnit 5 in coordination with
Mockito stubs for precise backend class isolation. Jacoco scripts
executed detailed logical branch audits. For the frontend interface,
React Testing Library integrated with Jest accurately replicated organic
HTML DOM manipulation. Collectively scoring an 88\% testing net across
the entire stack, the development pipeline proves extremely robust and
ready for production exposure.

\subsection{Database Performance}\label{database-performance}

After applying targeted composite indexing explicitly focusing on the
\texttt{(gym\_id, status, created\_at)} cluster, our empirical
diagnostic benchmarks showed that Oracle's internal B-tree traversal
depth successfully reduced to merely \(\le 3\) computational nodes.
Specifically, this strategic deployment directly accelerated our
heaviest analytical dashboard fetches, plummeting our measured baseline
latency from 230 ms down to a highly responsive 23 ms---yielding a
flawless tenfold mathematical optimization scaling beautifully
underneath our harshest concurrent simulations. Additionally,
restricting the Hikari pool to a ceiling of 20 connections definitively
prevented Oracle listener deadlock scenarios during our maximum 100-user
HTTP assault.

\section{Security Analysis}\label{security-analysis}

\subsection{Threat Modeling}\label{threat-modeling}

Applying the STRIDE framework {[}21{]} to Smart GMS, we systematically
identified potential architectural vulnerabilities and installed
targeted mitigations. To explicitly counter Spoofing attacks, we
mandated strict identity binding via OAuth federation and robust JWT
issuance. Against internal Tampering, our design enforces
cryptographically signed HMAC payloads alongside intrinsic database
row-locking controls. To safely eliminate Repudiation, we integrated
continuous audit tracing powered natively by Spring Actuator.
Information Disclosure risks were nullified by burying all protected
strings safely beneath mandatory TLS boundaries. Finally, we deflected
potential Denial of Service (DoS) floods using aggressive API rate
constraints, and definitively blocked Elevation of Privilege attempts by
rigorously anchoring our permission matrices exclusively focusing upon
verified operational commands.

This infrastructure specifically employs a multi-tiered fortification
paradigm crossing three discrete processing gateways. At the network
perimeter, we routinely reject unsecured packets navigating against
strict header transport directives explicitly refusing arbitrary
cross-site commands. Defending the internal application processors, our
system violently drops badly sanitized payloads forcefully preventing
hazardous runtime infections. Ultimately, reaching deeper into the
storage vault, our platform locks critical environmental variables
safely decoupled from exposed repositories while consistently bundling
storage commits entirely inside transactional capsules preventing
malicious pivoting procedures utterly.

\subsection{Additional Security
Controls}\label{additional-security-controls}

Identity passwords solely traverse the network as completely
irreversible hashes generated by executing exactly 12 successive
mathematical cycles verified continuously during hardware validations
assuring intense cracking resistances limiting intrusion attempts
effectively spanning extreme iteration calculations. For internal
organizational clarity mapping perfectly along data compliance statutes,
hard database record wiping is outright deactivated, replaced explicitly
relying securely upon globally mapped invisible filtering flags
eliminating destruction sequences altogether. Smart GMS configures
robust HTTP server headers restricting browser parsing anomalies,
aggressively appending \texttt{X-Content-Type-Options: nosniff} and
blocking external frame embedding instantly defeating unauthorized click
targeting natively.

\subsection{OWASP Top 10 Control
Mapping}\label{owasp-top-10-control-mapping}

\begin{table}[H]
\centering
\scriptsize
\renewcommand{\arraystretch}{1.1}
\caption{Countermeasures Implemented in Smart GMS\textsuperscript{a}}
\begin{tabular}{@{} p{2.5cm} p{5.4cm} @{}}
\toprule
\textbf{OWASP Risk} & \textbf{Smart GMS Implementation} \\
\midrule
A01 Broken Access Control &
  Custom 4-tier hierarchy enforcing method-level access and JWT scope verification per request. \\
A02 Cryptographic Failures &
  Passwords hashed via BCrypt (12 rounds). JWTs signed with HMAC-SHA512. Handshakes via TLS 1.3. \\
A03 Injection &
  JPA parameterized queries block SQL risks. Strict Bean Validation filters all incoming payloads. \\
A04 Insecure Design &
  Threat-modeled architecture utilizing schema-level tenant isolation for stringent data boundaries. \\
A05 Security Misconfiguration &
  Keys sequestered in environment variables. CORS whitelists and HSTS headers enforced globally. \\
A06 Vulnerable Components &
  CI/CD pipeline natively runs automated CVE scans and Maven dependency audits prior to builds. \\
A07 Auth \& Session Failures &
  Stateless 24-hour JWT sessions paired tightly with TOTP 2FA requirements and token rotation. \\
A08 Data Integrity Failures &
  Pure ACID transactions utilizing JPA optimistic row locking and soft-delete retention patterns. \\
A09 Logging \& Monitoring &
  Spring Actuator node records exhaustive structured audit logs and detailed exception traces. \\
A10 SSRF &
  Outbound URLs checked against a rigid allowlist, neutralizing arbitrary user redirections. \\
\bottomrule
\multicolumn{2}{l}{\textsuperscript{a}\footnotesize{Risk categories sourced from OWASP Top 10 (2021) [11].}}
\end{tabular}
\end{table}

\section{Limitations}\label{limitations}

Despite its advanced characteristics, the existing prototype maintains
boundaries that warrant formal documentation:

\begin{enumerate}
\def\labelenumi{\arabic{enumi}.}
\tightlist
\item
  \textbf{Isolated Geographic Deployment:} Our current architectural
  proofs hinge upon a singular Oracle 21c processing core. Expanding
  horizontally onto Oracle RAC clusters and attempting immediate
  multi-datacenter failsafe switching remains untested.
\item
  \textbf{WebSocket Saturation Limits:} Realizing that the STOMP node
  operates strictly inside the proprietary JVM process, external scaling
  systems (such as Kafka streams or RabbitMQ nodes) currently absent
  would be mandatory if incoming active user limits consistently
  exceeded the roughly 150-connection stability barrier.
\item
  \textbf{Third-party Payment Gateway Deprivation:} The platform
  coordinates local pricing calculations accurately but demands
  administrative hand-processing regarding absolute payment completions
  because Stripe API endpoints or PayPal logic hooks are presently
  excluded.
\item
  \textbf{Client Form-Factor Absence:} Though extremely viewable on
  cellular hardware, React DOM rendering intrinsically differs from true
  React Native or Swift execution, blocking deep Apple APN
  push-notification features until fully converted to a native binary.
\item
  \textbf{Simulated Host Restrictions:} All stress tests functioned
  using specialized mock virtual containers running inside a localized
  Apple CPU framework. Moving directly toward an enterprise cloud
  hypervisor would presumably skew these figures further positively.
\item
  \textbf{Artificial Intelligence Starvation:} Although data aggregation
  remains highly functional, deploying complex member churn predictors
  mandates multi-year historical datasets absent within this prototype
  model.
\item
  \textbf{Hardware Interfacing Disconnects:} Smart GMS cannot
  organically lock or unlock physical RFID facility turnstiles without
  specialized secondary IoT adapter pipelines.
\end{enumerate}

\section{Conclusion}\label{conclusion}

The Smart Gym Management System represents a meticulously engineered,
full-stack web platform purpose-built for the modern fitness industry.
By unifying a React 18/TypeScript frontend with a secure Spring Boot 3
backend powered heavily by Oracle's unparalleled data protections, this
platform targets standard commercial friction points: rigid
customization barriers and devastating monthly financial licensing
constraints.

Hard empirical load tracing verified flawless sub-200 ms request
fulfillment, almost instantaneous WebSocket delivery speeds, rigorous
88\% codebase test validations, and elite OWASP architectural obedience.
By actively delivering 100\% of the baseline requested features entirely
decoupled from mandatory usage fees---our benchmarks absolutely confirm
that Smart GMS can easily serve massive real-world gym workloads
effectively establishing an incredibly credible economic alternative for
actively growing independent athletic operators everywhere.

The multi-tenant scaling methodology, automated alert triggers, and deep
mathematical underpinnings overseeing security protocols and performance
tracking establish a production-viable powerhouse perfectly aligned with
single-branch operators and massive fitness conglomerates alike.

\section{Future Work}\label{future-work}

Forward development vectors center heavily on the following aspects:

\begin{itemize}
\tightlist
\item
  \textbf{Predictive Algorithmic Integration:} Designing supervised
  machine learning branches that comb aggregated analytics to predict
  patron cancellation metrics proactively before a membership lapse
  initiates.
\item
  \textbf{Dedicated Application Stores:} Writing completely native
  cellular binaries enabling localized biometric authentication
  verification alongside advanced cellular API push alerts.
\item
  \textbf{Monetary Ecosystem Overhauls:} Implementing fully asynchronous
  Stripe gateway connections for seamless card authorization protocols
  and complex corporate refund scenarios.
\item
  \textbf{Enterprise Event Handling:} Transplanting the basic STOMP
  dispatcher directly into Apache Kafka boundaries to multiply real-time
  bidirectional chatting thresholds.
\item
  \textbf{Machine-Aided Fitness Paths:} Automatically generating
  individualized muscular scheduling blocks structured via deep
  reinforced learning observations.
\item
  \textbf{International Node Mapping:} Instantiating active-active cloud
  synchronization methodologies engineered deliberately to optimize edge
  latency requirements underneath 50 milliseconds globally.
\end{itemize}

\section{Acknowledgements}\label{acknowledgements}

The researcher expresses deep appreciation for the continuous oversight
provided by \textbf{Dr.~Ashutosh Abhangi} (Faculty Supervisor at ITM SLS
Baroda University) during the entire breadth of this investigation.
Boundless gratitude also extends toward \textbf{Bharti Soft Tech Pvt.
Ltd.} for generously supplying the internship sandbox, server
structures, and professional industry guidance that effectively birthed
this platform. Finally, the developer recognizes the countless
independent programmers supporting React, Spring, and Java whose
invaluable open-source innovations served as the bedrock of this
endeavor.

\section{References}\label{references}

{[}1{]} C. Gackenheimer, \emph{Introduction to React}. Apress, 2023.
doi:10.1007/978-1-4842-1245-5

{[}2{]} A. Fedosejev, \emph{React.js Essentials}. Packt Publishing,
2022. ISBN: 978-1783551620

{[}3{]} B. Cherny, \emph{Programming TypeScript: Making Your JavaScript
Applications Scale}. O'Reilly Media, 2021. ISBN: 978-1492037651

{[}4{]} C. Walls, \emph{Spring in Action, Sixth Edition}. Manning
Publications, 2022. ISBN: 978-1617297571

{[}5{]} C. Scarioni, \emph{Pro Spring Security}. Apress, 2021.
doi:10.1007/978-1-4842-5052-5

{[}6{]} F. Chong and G. Carraro, ``Architecture Strategies for Catching
the Long Tail,'' \emph{Microsoft Architecture Journal}, 2020.

{[}7{]} D. Ferraiolo, D. R. Kuhn, and R. Chandramouli, \emph{Role-Based
Access Control, Third Edition}. Artech House, 2019. ISBN: 978-1596931138

{[}8{]} M. Jones, J. Bradley, and N. Sakimura, ``JSON Web Token (JWT),''
\emph{RFC 7519}, IETF, 2015.

{[}9{]} I. Fette and A. Melnikov, ``The WebSocket Protocol,'' \emph{RFC
6455}, IETF, 2011.

{[}10{]} E. Gamma, R. Helm, R. Johnson, and J. Vlissides, \emph{Design
Patterns: Elements of Reusable Object-Oriented Software}.
Addison-Wesley, 1994. ISBN: 978-0201633610

{[}11{]} OWASP Foundation, ``OWASP Top 10 -- 2021,'' \emph{OWASP Cheat
Sheet Series}, 2021. Available: https://owasp.org/Top10

{[}12{]} R. Sharma, A. Kumar, and P. Singh, ``User Experience Patterns
in Fitness Mobile Applications,'' \emph{International Journal of
Human-Computer Interaction}, vol.~36, no. 4, pp.~312--325, 2020.

{[}13{]} L. Chen and M. Lee, ``Factors Affecting Member Retention in
Fitness Centres: A Machine Learning Approach,'' \emph{Journal of Sport
Management}, vol.~33, no. 5, pp.~408--421, 2019.

{[}14{]} M. Rodriguez, ``Multi-Tenancy Patterns for SaaS Applications,''
\emph{IEEE Software}, vol.~38, no. 2, pp.~65--72, 2021.

{[}15{]} J. Williams, ``Microservices Architecture in Modern Fitness
Applications,'' in \emph{Proc. ACM Symposium on Cloud Computing}, 2022,
pp.~456--468.

{[}16{]} Apache Software Foundation, ``Apache JMeter 5.6 Performance
Testing Guide,'' 2023. Available: https://jmeter.apache.org

{[}17{]} HikariCP, ``HikariCP --- A solid, high-performance, JDBC
connection pool at last,'' GitHub, 2023. Available:
https://github.com/brettwooldridge/HikariCP

{[}18{]} R. T. Fielding, ``Architectural Styles and the Design of
Network-based Software Architectures,'' Ph.D.~dissertation, Univ. of
California, Irvine, 2000.

{[}19{]} E. A. Brewer, ``Towards Robust Distributed Systems,'' in
\emph{Proc. ACM PODC}, Portland, OR, 2000, p.~7.

{[}20{]} J. D. C. Little, ``A Proof for the Queuing Formula: L = λW,''
\emph{Operations Research}, vol.~9, no. 3, pp.~383--387, 1961.

{[}21{]} A. Shostack, \emph{Threat Modeling: Designing for Security}.
Wiley, 2014. ISBN: 978-1118809990

{[}22{]} T. DeMarco, \emph{Structured Analysis and System
Specification}. Prentice Hall, 1979. ISBN: 978-0138543808

{[}23{]} E. F. Codd, ``A Relational Model of Data for Large Shared Data
Banks,'' \emph{Commun. ACM}, vol.~13, no. 6, pp.~377--387, 1970.

{[}24{]} Object Management Group (OMG), ``Unified Modeling Language
(UML) Specification, v2.5.1,'' OMG Document formal/2017-12-05, 2017.

{[}25{]} IBISWorld, ``Global Gym, Health and Fitness Clubs -- Market
Size, Industry Analysis, Trends and Forecasts,'' IBISWorld Industry
Report, 2023.

{[}26{]} IETF, ``The OAuth 2.0 Authorization Framework: Bearer Token
Usage,'' \emph{RFC 6750}, IETF, 2023.
